% =====================================================================
%  CONJURA CONJECTURE STATEMENT
%  No expanding weak pseudorandom generators computed by quadratic
%  polynomials over the integers.
%  Requires conjura-conjecture.cls in the same directory.
% =====================================================================
\documentclass{conjura-conjecture}

\runninghead{NO EXPANDING WEAK QUADRATIC PRGS}

\newcommand{\ZZ}{\mathbb{Z}}
\newcommand{\RR}{\mathbb{R}}
\newcommand{\NN}{\mathbb{N}}
\newcommand{\EE}{\mathbb{E}}
\newcommand{\cA}{\mathcal{A}}
\newcommand{\cX}{\mathcal{X}}
\newcommand{\cQ}{\mathcal{Q}}
\newcommand{\Adv}{\mathrm{Adv}}

\begin{document}

\cjkicker{CONJURA \textperiodcentered{} OPEN PROBLEM}
\cjtitle{No Expanding Weak Quadratic Pseudorandom Generators}
\cjsubtitle{Arbitrary integer quadratic maps at stretch
  $n^{1+\varepsilon}$}
\cjstatus{Statement: AI-written from Boaz Barak, Samuel B. Hopkins,
  Aayush Jain, Pravesh Kothari, Amit Sahai, \emph{Sum-of-Squares Meets
  Program Obfuscation, Revisited}, IACR Cryptology ePrint Archive 2018,
  not yet checked by a human. Settled only in the restricted case where
  $q_1, \dots, q_m$ are drawn independently and identically from a
  \emph{nice} distribution --- supported on homogeneous (no linear
  term) degree-two polynomials, with
  $\|q\|_2^2 \le C \cdot \EE\|q\|_2^2$ for a constant $C$, with
  $\mathrm{Var}(Q_{i,j}) = 1$ and with the coefficients $Q_{i,j}$
  pairwise independent --- and where the input is polynomially bounded
  and sub-Gaussian: there the paper's Theorem~2 gives a polynomial-time
  algorithm that recovers $x$ exactly from $m \ge n(\log n)^{O(1)}$
  samples, which in particular distinguishes. The hypothesis as stated,
  with arbitrary (possibly correlated) polynomials and an arbitrary
  bounded input distribution, is open.}
\cjcategory{\emph{Public-Key Cryptography \& Assumptions}.}

\vspace{0.4em}

\begin{informalconjecture}
Several candidate constructions of indistinguishability obfuscation
reduce security to a very weak kind of generator: you publish a list of
quadratic polynomials with small integer coefficients, evaluate all of
them at one secret input drawn from some distribution, and you only
need the resulting list of values to be indistinguishable from that
same list with small independent integer noise added to each entry.
Nothing needs to look uniformly random, and the secret need not be
recoverable even in principle. This paper conjectures that no such
object exists once the number of polynomials exceeds the number of
variables by any fixed polynomial factor: whenever you have more than
about $n$ polynomials in $n$ variables, all with polynomially bounded
integer coefficients, and the input distribution is also polynomially
bounded, there is an efficient distinguisher for that family which
achieves a constant advantage, the constant depending only on the
stretch and the coefficient bound. The authors prove this in the
special case where the polynomials are drawn independently from a
distribution on homogeneous degree-two polynomials satisfying mild
normalisation and pairwise-independence conditions, using a
semidefinite program, and they report that they know of no degree-two
candidate their attacks do not break. The general case, where the
polynomials may be correlated with each other or chosen adversarially,
is left open.
\end{informalconjecture}

\section{The setting}\label{sec:setting}

Candidate constructions of indistinguishability obfuscation (iO) have
converged on a common shape: bootstrap iO from a low-degree
cryptographic multilinear map, LWE, and a low-degree expanding
generator over the integers. Lin and Tessaro~\cite{LT} reduced iO to
bilinear maps together with a block-wise two-local pseudorandom
generator of superlinear stretch. Barak, Brakerski, Komargodski and
Kothari~\cite{BBKK}, and independently Lombardi and
Vaikuntanathan~\cite{LV}, then showed that the Lin--Tessaro candidate,
and indeed any generator with their required parameters, can be broken
using the degree-two sum-of-squares semidefinite program.

The works of Ananth, Jain and Sahai~\cite{AJS}, Agrawal~\cite{Agr18}
and Lin and Matt~\cite{LM} proposed a way around that obstacle by
weakening the requirement in two directions at once. Structurally, the
outputs need only be \emph{degree-two polynomials} of the input, rather
than block-wise local. Security-wise, one asks only that the vector of
evaluations $q_1(x), \dots, q_m(x)$ be indistinguishable from the same
vector with small independent integer perturbations added --- the
\emph{$\Delta$RG} property of~\cite{AJS}, and the closely related
\emph{pseudo flawed-smudging} property of~\cite{LM}. Nothing is
required to look uniform, and the input need not be recoverable.

This paper conjectures that the retreat does not help: at stretch
$n^{1+\varepsilon}$ no such object exists, for \emph{any} choice of
quadratic polynomials with polynomially bounded integer coefficients
and \emph{any} polynomially bounded input distribution. The conjecture
is stated as a hypothesis about the existence of an efficient
distinguisher for each such family, so a proof is an algorithm and a
refutation is a construction.

Two obstructions are recorded. First, the obvious route to a
distinguisher --- recover $x$ from the evaluations --- is not available
in general: it need not even be information-theoretically possible to
recover $x$, and where it is, doing so is an instance of the NP-hard
problem of solving quadratic equations. Second, the paper's own theorem
covers only the case where the polynomials are drawn
\emph{independently and identically} from a distribution supported on
homogeneous degree-two polynomials and satisfying normalisation and
pairwise-independence conditions (called \emph{nice}); the conjecture
allows arbitrary, possibly heavily correlated, polynomials.

The technical route to the theorem goes through low-rank recovery.
Writing a homogeneous quadratic as
$q_i(x) = \langle Q_i, xx^{\top} \rangle$ turns seed recovery into
recovering a rank-one matrix from $m$ linear measurements, which
nuclear-norm (trace-norm) minimisation provably solves when the
measurement matrices satisfy matrix RIP~\cite{RFP} or the incoherence
condition of Gross~\cite{Gross}. Both conditions are randomness
conditions on the measurements; a proof of the conjecture would have to
dispense with them entirely, since the polynomials are adversarial.

Progress to date. Theorem~2 of the paper proves the i.i.d.-nice case: a
polynomial-time sum-of-squares/trace-norm algorithm recovers $x$
exactly with probability $1 - n^{-\log n}$ from
$m \ge n(\log n)^{O(1)}$ observations, for every nice distribution
$\cQ$ --- that is, every distribution supported on homogeneous (no
linear term) degree-two polynomials with
$\|q\|_2^2 \le C \cdot \EE\|q\|_2^2$ for a constant $C$, with
$\mathrm{Var}(Q_{i,j}) = 1$ and with the coefficients $Q_{i,j}$
pairwise independent --- and every input uniform on
$\{-t, \dots, t\}^n$ with $t \le n^{O(1)}$ (and more generally for
centred $O(\EE\|x\|^2/n)$-sub-Gaussian real inputs, whose coordinates
need not be independent). This covers random dense and random sparse
polynomials, and hence breaks the Ananth--Jain--Sahai and Agrawal
candidates. The paper additionally reports experiments that break every
degree-two candidate the authors are aware of, including the Lin--Matt
one, and states that it knows of no degree-two construction that
escapes a variant of the attack.

\section{Notation and parameters}\label{sec:notation}

\begin{itemize}[leftmargin=1.2em]
  \item $n \in \NN$ is the number of variables; all asymptotics are as
  $n \to \infty$, and $[m] := \{1, \dots, m\}$.
  \item A \emph{quadratic polynomial} is a polynomial
  $q : \RR^n \to \RR$ of total degree at most $2$. Such a polynomial is
  presented to an algorithm as its vector of coefficients.
  \item A function $\Lambda : \NN \to \NN$ is \emph{polynomially
  bounded} if $\Lambda(n) \le n^{O(1)}$.
  \item An algorithm $\cA$ is \emph{efficient} if it is probabilistic
  and runs in time polynomial in the bit length of its input.
  \item For distributions $D_1, D_2$ over a common domain, the
  \emph{advantage} of an algorithm $\cA$ is
  \begin{align*}
    \Adv_{\cA}(D_1, D_2) \;:=\; \Bigl|
      &\Pr_{z \sim D_1}[\cA(z) = 1] \\
      &{}- \Pr_{z \sim D_2}[\cA(z) = 1] \Bigr| .
  \end{align*}
\end{itemize}

The parameters of the statement below are as follows.

\begin{itemize}[leftmargin=1.2em]
\item $n$: number of variables of each polynomial, i.e.\ the length of
  the secret input.
\item $m$: number of polynomials published, i.e.\ the output length of
  the generator. No upper bound is imposed; efficiency is measured in
  the input length, which grows with $m$.
\item $\varepsilon$: stretch exponent; the conjecture assumes
  $m \ge n^{1+\varepsilon}$.
\item $\Lambda(n)$: polynomial bound on the coefficients of the
  polynomials, on the support of the input distribution, and on the
  support of the noise.
\item $\cX$: distribution of the secret input over $\ZZ^n$.
\item $\Delta_i$: distribution of the noise added to the $i$-th output;
  required to have no atom of mass $0.9$ or more.
\item $c$: the constant distinguishing advantage the conjectured
  algorithm achieves.
\end{itemize}

\section{The conjecture}\label{sec:conjectures}

\begin{definition}[$\Lambda$-bounded polynomial]\label{def:boundedpoly}
An $n$-variate polynomial $q$ is \emph{$\Lambda$-bounded} if all of
$q$'s coefficients are integers in the interval
$[-\Lambda, +\Lambda]$.
\end{definition}

\begin{definition}[$\Lambda$-bounded distribution]\label{def:bounddist}
A distribution $\cX$ over $\ZZ^n$ is \emph{$\Lambda$-bounded} if it is
supported on $[-\Lambda, +\Lambda]^n$. A distribution $\Delta$ over
$\ZZ$ is \emph{$\Lambda$-bounded} if it is supported on
$[-\Lambda, +\Lambda]$.
\end{definition}

\begin{conjecture}[no expanding weak quadratic
  generators]\label{conj:main}
For every constant $\varepsilon > 0$ and every polynomially bounded
$\Lambda : \NN \to \NN$ there is a constant $c > 0$ such that the
following holds for all sufficiently large $n \in \NN$.

Let $m \ge n^{1+\varepsilon}$ be an integer, let
$q_1, \dots, q_m : \RR^n \to \RR$ be quadratic $\Lambda(n)$-bounded
polynomials, let $\cX$ be a $\Lambda(n)$-bounded distribution over
$\ZZ^n$, and for each $i \in [m]$ let $\Delta_i$ be a
$\Lambda(n)$-bounded distribution over $\ZZ$ such that
$\Pr[\Delta_i = z] < 0.9$ for every $z \in \ZZ$. Define the two
distributions
\begin{align*}
  D_1 &\;:=\; \bigl( q_1, \dots, q_m,\; q_1(x), \dots, q_m(x) \bigr), \\
  D_2 &\;:=\; \bigl( q_1, \dots, q_m, \\
      &\qquad\quad\; q_1(x) + \delta_1, \dots, q_m(x) + \delta_m \bigr),
\end{align*}
where $x \sim \cX$ in both, and where in $D_2$ the perturbations
$\delta_i \sim \Delta_i$ are drawn independently of one another and of
$x$. Then there exists an efficient algorithm $\cA$ with
\[
  \Adv_{\cA}(D_1, D_2) \;\ge\; c .
\]

Equivalently: there is no family of quadratic $\Lambda(n)$-bounded
polynomials, of stretch $m \ge n^{1+\varepsilon}$, whose evaluations at
a $\Lambda(n)$-bounded random input are computationally
indistinguishable from those evaluations perturbed by bounded
independent noise no one value of which is taken with probability $0.9$
or more.
\end{conjecture}

\begin{remark}[readings fixed by this write-up]\label{rem:readings}
Two choices above are the page author's rather than the paper's. First,
the paper writes ``then there exists an algorithm $\cA$ that can
distinguish'' without the word \emph{efficient}; polynomial time has
been made explicit here, because the very next paragraph of the paper
reasons about ``an efficient algorithm to recover $x$'', and because
the accompanying Theorem~2 supplies a polynomial-time algorithm. Read
literally without an efficiency requirement the hypothesis is far
weaker, so this is almost certainly the intended reading --- but it is
a reading, not a transcription. Second, the paper imposes no upper
bound on $m$ and none is imposed above; since efficiency is measured in
the bit length of the input, which contains all $m$ polynomials, the
running time allowed to $\cA$ grows with $m$. Note also that, following
the paper's quantifier order, $\cA$ is quantified after the instance
and so may depend on $m$, on $q_1, \dots, q_m$, on $\cX$ and on the
$\Delta_i$; only the advantage $c$ is uniform over instances.
\end{remark}

\begin{remark}[scope of the quantification --- the page author's
  commentary, postdating the paper]\label{rem:scope}
The following is not drawn from the source and postdates it. The
hypothesis fixes the polynomials independently of the input
distribution: the same $q_1, \dots, q_m$ appear in both distributions
and are not correlated with $x$. Later constructions of iO which use
degree-two maps evade attacks of this kind by sampling public data
jointly with the secret seed, so that the effective polynomials depend
on the secret~\cite{JLS}; such objects fall outside the quantification
above and are not refutations of it. This reading of the later
literature has not been verified, and a reviewer should confirm that no
refutation has appeared.
\end{remark}

\section{Bibliography}

\begin{conjurabibliography}{99}

\bibitem[Agr18]{Agr18}
S. Agrawal.
\emph{New methods for indistinguishability obfuscation: bootstrapping
and instantiation}.
IACR Cryptology ePrint Archive, 2018:633, 2018.

\bibitem[AJS]{AJS}
P. Ananth, A. Jain, A. Sahai.
\emph{Indistinguishability obfuscation without multilinear maps: iO
from LWE, bilinear maps, and weak pseudorandomness}.
IACR Cryptology ePrint Archive, 2018:615, 2018.

\bibitem[BBKK]{BBKK}
B. Barak, Z. Brakerski, I. Komargodski, P. Kothari.
\emph{Limits on low-degree pseudorandom generators (or: sum-of-squares
meets program obfuscation)}.
Electronic Colloquium on Computational Complexity (ECCC), 24:60, 2017.

\bibitem[Gross]{Gross}
D. Gross.
\emph{Recovering low-rank matrices from few coefficients in any basis}.
IEEE Transactions on Information Theory, 57(3):1548--1566, 2011.

\bibitem[JLS]{JLS}
A. Jain, H. Lin, A. Sahai.
\emph{Indistinguishability obfuscation from well-founded assumptions}.
STOC, 2021 (recalled from memory; not in this paper's reference list).
[UNVERIFIED]

\bibitem[LM]{LM}
H. Lin, C. Matt.
\emph{Pseudo flawed-smudging generators and their application to
indistinguishability obfuscation}.
IACR Cryptology ePrint Archive, 2018:646, 2018.

\bibitem[LT]{LT}
H. Lin, S. Tessaro.
\emph{Indistinguishability obfuscation from bilinear maps and
block-wise local PRGs}.
Cryptology ePrint Archive, Report 2017/250, 2017.

\bibitem[LV]{LV}
A. Lombardi, V. Vaikuntanathan.
\emph{On the non-existence of blockwise 2-local PRGs with applications
to indistinguishability obfuscation}.
IACR Cryptology ePrint Archive, 2017:301, 2017.

\bibitem[RFP]{RFP}
B. Recht, M. Fazel, P. A. Parrilo.
\emph{Guaranteed minimum-rank solutions of linear matrix equations via
nuclear norm minimization}.
SIAM Review, 52(3):471--501, 2010.

\end{conjurabibliography}

\end{document}
